\section{Periodic lattice memory on the torus}
\label{sec:lattice}

We now store successive configurations in a single torus coordinate.
The earlier records remain in that coordinate, but a residue calculation
removes them when the next transition is read.  A second coordinate
records whether the current state is halting.  Both the new record and
the halting signal become much smaller at each step.  This separation of
scales will make every derivative of the velocity and force decay faster
than every inverse power of time. The positional stack operations belong
to the generalized-shift approach \cite{Moore}; the residue arithmetic
and smooth realization are proved below.

The planar motion used for this purpose need not preserve area.  We
first construct and verify that motion.  A third component will then
cancel its divergence while leaving the coding plane invariant.

\begin{theorem}[Periodic lattice memory]\label{thm:lattice}
Fix a computable viscosity $\nu>0$.  A finite deterministic machine and
a finite input effectively determine a smooth force $g$ on the unit
flat torus $\T^3=(\R/\mathbb Z)^3$ with zero spatial mean.  The
Navier--Stokes equations with force $g$ and zero initial velocity have
a smooth solution $(U,0)$ with zero spatial mean.  Its particle from
\[
 a_*=(1/4,1/2,1/2)
\]
enters the fixed open set
\[
 O=\{x\in\T^3:1/2<x_1<1\pmod 1\}
\]
at some finite time if and only if the machine halts on that input.
For every pair of nonnegative integers $r,J$ and every spatial
multi-index $\alpha$, there is an effective constant $C_{r,\alpha,J}$
such that
\begin{equation}\label{eq:lattice-rapid}
 \|\partial_t^r\partial_x^\alpha U(t)\|_\infty
 +\|\partial_t^r\partial_x^\alpha g(t)\|_\infty
 \le C_{r,\alpha,J}(1+t)^{-J},\qquad t\ge0.
\end{equation}
There is also an effective smooth mean-zero solenoidal force $f$
with the same bounds and exactly the same velocity and material event.
It is $f=g-\nabla\phi$, with pressure $-\phi$, where $\phi$ is the
mean-zero periodic solution of $\Delta\phi=\operatorname{div}g$.

For either force, uniqueness holds among periodic classical solutions
for which the velocity, its time derivative and its spatial derivatives
through order two, and the pressure and its spatial gradient, are
continuous on every $[0,T]\times\T^3$.  Pressure is normalized to have
zero spatial mean.
\end{theorem}

\subsection{Reading a configuration from an integer residue}

Normalize a missing instruction by entering a fresh halting state without
moving or changing the scanned symbol. To retain absolute
tape position as well as the contents relative to the head, add a
read-only bit track with a unique marked bit at the initial origin.
Transitions preserve the bit and ignore it when choosing an instruction.
The resulting alphabet is denoted by $\Gamma$, its all-blank symbol by
$\square$, its state set by $Q$, and its finite initial word by $\omega$.
The word contains the marked origin even when the original input is
empty.  This modification preserves halting.

Choose $b=\max(2,|\Gamma|,|Q|)$.  Give the symbols distinct digits in
$\{0,\ldots,b-1\}$, assigning zero to $\square$, and number the states
from zero.  After $n$ transitions, put
\[
 d_n=|\omega|+1+2n,\qquad p_n=2d_n+1,\qquad P_n=b^{p_n}.
\]
The nearest $d_n$ cells to the left of the head form a stack with
integer code $L$; the nearest $d_n$ cells starting at the head form a
stack with code $R$.  In each stack the nearest cell is the least
significant base-$b$ digit.  All cells outside these two stacks are
blank: a head travels at most one cell per transition, whereas the
available width grows by two.  If the state number is $s$, the complete
configuration therefore has the code
\begin{equation}\label{eq:lattice-code}
 c_n=L+b^{d_n}R+b^{2d_n}s,\qquad 0\le c_n<P_n.
\end{equation}
The slot number $n$ specifies the stack width.  The origin bit determines
the head's absolute position, so the code loses no tape information.

We must prescribe the field at every spatial address, rather than only
at addresses visited by this computation.  For an arbitrary integer
$j$, let $[j]_q$ denote its least nonnegative residue modulo the positive
integer $q$.  Extract
\[
 \begin{split}
 c&=[j]_{P_n},\\
 L&=[c]_{b^{d_n}},\qquad
 R=\bigl[\lfloor c/b^{d_n}\rfloor\bigr]_{b^{d_n}},\\
 s&=\lfloor c/b^{2d_n}\rfloor,\qquad
 l=[L]_b,\quad a=[R]_b.
 \end{split}
\]
Define $h_n(j)=1$ if $s$ is a halting-state index, and $h_n(j)=0$
otherwise.  If $s$ is a valid nonhalting state and $a$ a valid symbol,
look up the single instruction $(s,a)\mapsto(s',a',m)$, where
$m\in\{-1,0,1\}$.  Set
\begin{equation}\label{eq:lattice-transition}
 \begin{split}
 (L',R')&=
 \begin{cases}
 (a'+bL,(R-a)/b),&m=1,\\
 ((L-l)/b,l+ba'+b(R-a)),&m=-1,\\
 (L,R-a+a'),&m=0,
 \end{cases}\\
 F_n(j)&=L'+b^{d_{n+1}}R'+b^{2d_{n+1}}s'.
 \end{split}
\end{equation}
In every other case define $F_n(j)=0$.  Digits away from the scanned
cell need not be checked for membership in the alphabet.

\begin{lemma}\label{lem:lattice-address}
The functions $F_n,h_n$ on $\mathbb Z$ are $P_n$-periodic,
$h_n\in\{0,1\}$, and $0\le F_n<P_{n+1}$ at every address.
If $c_n$ is a true nonhalting configuration and
$j\equiv c_n\pmod {P_n}$, then $F_n(j)=c_{n+1}$.
\end{lemma}
\begin{proof}
All extracted data depend only on $[j]_{P_n}$.  In each nonzero branch,
$0\le L',R'<b^{d_n+1}$.  For the left move, the only less immediate
bound follows from
\[
 R-a\le b^{d_n}-b,\qquad l,a'\le b-1,
\]
which give $l+ba'+b(R-a)\le b^{d_n+1}-1$.  The right and stationary
branches satisfy the same bounds directly.  Since $d_{n+1}=d_n+2$
and $0\le s'<b$, the output fits in the three fields of the next code.
The formulas are precisely the stack operations for a tape transition.
In particular, after a left move the new right stack consists of the
old left top, the symbol just written, and the old right tail.  This is
$l+ba'+b(R-a)$; the last term has two digit shifts because $R-a$
already contains one factor of $b$.  The two new blank places suffice
for every branch.
\end{proof}

\subsection{Scales that retain old records without reading them}

Suppose a real coordinate stores a partial history as
$y_*+\sum_{k=0}^n\epsilon_kc_k$, with $y_*=1/2$.
To recover $c_n$, we want every earlier summand, after division by
$\epsilon_n$, to be an integer multiple of $P_n$.  We also want the
next displacement to be much smaller than $\epsilon_n$, so a smooth
selector will read the same integer throughout that displacement.
The following explicit scales meet both requirements:
\begin{equation}\label{eq:lattice-scales}
 K_n=(p_0+10)2^{n(n+1)},\qquad
 \epsilon_n=b^{-K_n},\qquad
 \beta_0=1/4,\quad \beta_n=\epsilon_n\quad(n\ge1).
\end{equation}
The numbers $\beta_n$ give the distance of a nonhalting signal from
the boundary $x_1=1/2$. The separate choice $\beta_0=1/4$ fixes the
initial first coordinate at $1/4$, independently of the input; later
distances use the shrinking memory scales. Its contribution to the
first transition is a single finite-time term in the decay estimates.

\begin{lemma}\label{lem:lattice-scales}
For $n\ge0$,
\begin{align}
 P_n&\mid\epsilon_n^{-1},&
 P_n&\mid\epsilon_k/\epsilon_n\quad(0\le k<n),
 \label{eq:lattice-divisibility}\\
 P_{n+1}\epsilon_{n+1}
 &\le\epsilon_n^{\,2^{2n+2}-1}
 \le\epsilon_n^3<\epsilon_n/4.
 \label{eq:lattice-plateau-bound}
\end{align}
Moreover, $\beta_n$ decreases strictly and $\epsilon_n\le2^{-n}$.
\end{lemma}
\begin{proof}
Here $p_n=p_0+4n$ and $K_{n+1}=2^{2n+2}K_n$.
Induction gives $K_n\ge p_{n+1}$, starting with
$K_0=p_0+10\ge p_0+4$.  Thus $K_n\ge p_n$ and, for $k<n$,
\[
 K_n-K_k\ge K_n-K_{n-1}\ge3K_{n-1}\ge p_n.
\]
These inequalities prove the two divisibilities, since all quantities
are powers of the same integer $b$.  Also
\[
 b^{p_{n+1}-K_{n+1}}
 \le b^{K_n-2^{2n+2}K_n}
 =\epsilon_n^{\,2^{2n+2}-1}.
\]
The remaining assertions follow from $b\ge2$, $K_0\ge13$ and the
strict growth of $K_n$.
\end{proof}

We have now specified the finite rule table and the scale arithmetic.
The next step makes the residue calculation into a smooth field whose
actual particle follows the whole stored history.

\subsection{Smooth selectors and the complete particle path}

Fix an effective smooth nondecreasing function $\sigma$ with
$\sigma(t)=0$ for $t\le1/3$ and $\sigma(t)=1$ for $t\ge2/3$.
Take $\sigma(t)=H(3t-1)$ with the function $H$ fixed in
Section~\ref{sec:analytic}.
Choose an effective smooth bump $\psi$ that is one on $[-1/4,1/4]$
and zero outside $[-3/8,3/8]$.  Products of translates and rescalings
of $H$ supply such a bump.  These functions and their derivatives have
effective bounds; flatness at the endpoints follows from the
exponential factor in $\rho$.

For $j\in\mathbb Z$, define a bump centered at the $j$th point of
the scale-$n$ lattice by
\[
 J_{n,j}(y)=\psi\bigl((y-y_*)/\epsilon_n-j\bigr).
\]
Their supports are pairwise disjoint.  In planar coordinates
$\xi=(\xi_1,\xi_2)$, prescribe the loading field and the $n$th
transition field by
\begin{align}
 v_{\rm in}(t,\xi)&=\sigma'(t)(0,\epsilon_0c_0),
 \label{eq:lattice-loader}\\
 v_n(t,\xi)&=\sigma'(t-1-n)
 \begin{pmatrix}
  \beta_n-\beta_{n+1}
     +2\beta_{n+1}\displaystyle\sum_{j\in\mathbb Z}h_n(j)J_{n,j}(\xi_2)\\
  \epsilon_{n+1}\displaystyle\sum_{j\in\mathbb Z}F_n(j)J_{n,j}(\xi_2)
 \end{pmatrix},
 \label{eq:lattice-slot}\\
 v&=v_{\rm in}+\sum_{n\ge0}v_n.
 \label{eq:lattice-planar-field}
\end{align}
The field is independent of $\xi_1$.  Translation of $\xi_2$ by
one shifts the lattice index by $\epsilon_n^{-1}$, a multiple of
$P_n$ by Lemma~\ref{lem:lattice-scales}.  The coefficients therefore
repeat, and $v$ is a field on $\T^2$.  Only the loader acts during
$[0,1]$, and only $v_n$ acts during $[1+n,2+n]$.  The fixed zero
collars of the pulses make their sum smooth, including at every slot
boundary and at time zero.

Let $Y(t)$ be the real-coordinate lift of its particle starting at
$(1/4,y_*)$.  The loader leaves the first coordinate unchanged and
ends at
\[
 Y(1)=(1/4,y_*+\epsilon_0c_0).
\]
We claim that, at $t_n=1+n$ and up to the first halting configuration,
\begin{equation}\label{eq:lattice-history}
 Y_1(t_n)=1/2-\beta_n,\qquad
 Y_2(t_n)=y_*+\sum_{k=0}^n\epsilon_kc_k.
\end{equation}
The claim holds for $n=0$.  Assuming it at $n$, put
\[
 j_*=(Y_2(t_n)-y_*)/\epsilon_n.
\]
This is an integer, and \eqref{eq:lattice-divisibility} gives
$j_*\equiv c_n\pmod {P_n}$.  Changing the lift of $Y_2$ by an
integer changes $j_*$ by a multiple of $P_n$, so the decoded instruction
is intrinsic to the torus.

Throughout this slot the actual trajectory is
\begin{align}
 Y_2(t_n+\tau)&=Y_2(t_n)+\sigma(\tau)\epsilon_{n+1}F_n(j_*),
 \label{eq:lattice-path-y}\\
 Y_1(t_n+\tau)&=Y_1(t_n)+\sigma(\tau)
  \bigl(\beta_n-\beta_{n+1}+2\beta_{n+1}h_n(j_*)\bigr),
 \quad 0\le\tau\le1.
 \label{eq:lattice-path-x}
\end{align}
Indeed, Lemma~\ref{lem:lattice-address} and
\eqref{eq:lattice-plateau-bound} bound the second-coordinate
displacement by $\epsilon_n/4$.  The proposed curve stays on the
plateau $J_{n,j_*}=1$, where every other bump vanishes.  Differentiating
\eqref{eq:lattice-path-y}--\eqref{eq:lattice-path-x} therefore gives
exactly \eqref{eq:lattice-slot}.  Spatial smoothness and ODE uniqueness
identify it with the actual particle.

At a nonhalting configuration, Lemma~\ref{lem:lattice-address} gives
$F_n(j_*)=c_{n+1}$ and $h_n(j_*)=0$.  The endpoint is thus the next
history in \eqref{eq:lattice-history}.  Throughout that slot the first
coordinate lies between $1/2-\beta_n$ and $1/2-\beta_{n+1}$, both
strictly below $1/2$.  At a halting configuration $h_n(j_*)=1$, and
the first-coordinate endpoint is instead
\[
 Y_1(t_n+1)=1/2+\beta_{n+1}\in(1/2,1).
\]
This also handles an initially halted input, whose signal appears in
the first slot after loading.  There is no first-coordinate wrap before
the first detection.  Consequently the open event has the required
equivalence at all real times, not just at integer endpoints.

\subsection{Rapid bounds and the incompressible lift}

The stored history grows, but neither its length nor the number of
lattice cells enters a derivative bound.  At each spatial point at most
one bump in \eqref{eq:lattice-slot} is nonzero.  For $|\alpha|=k$,
all time orders $r$, and $n\ge0$, the coefficient bounds give
\begin{equation}\label{eq:lattice-slot-bound}
 \|\partial_t^r\partial_\xi^\alpha v_n\|_\infty
 \le C_{r,k}\left(\mathbf 1_{k=0}\beta_n
       +P_{n+1}\epsilon_{n+1}\epsilon_n^{-k}\right).
\end{equation}
The constant drift disappears after a spatial derivative; each derivative
of a bump costs one factor of $\epsilon_n^{-1}$.  For
$n\ge\max(1,k)$, \eqref{eq:lattice-plateau-bound} bounds the
right-hand side by $C'_{r,k}\epsilon_n$.  Since
$\epsilon_n\le2^{-n}$ and $1+t\le n+3$ on the $n$th slot,
every polynomial time weight is bounded.  The finitely many earlier
slots and the loader contribute finite effective constants.  Thus $v$
satisfies all the mixed rapid bounds in \eqref{eq:lattice-rapid}.

Choose an effective smooth periodic function $\eta$ on $\T$ with
$\eta(1/2)=0$ and $\eta'(1/2)=1$.  One explicit choice is
$(z-1/2)$ times a smooth cutoff supported in
$(1/2-3/16,1/2+3/16)$ and equal to one near $1/2$, extended periodically.
Define on $\T^3$
\begin{equation}\label{eq:lattice-lift}
 U(t,\xi,z)=
 \bigl(\eta'(z)v_1(t,\xi),\eta'(z)v_2(t,\xi),
       -\eta(z)\operatorname{div}_\xi v(t,\xi)\bigr).
\end{equation}
The two divergence terms cancel.  The first two components have zero
mean because $\int_\T\eta'=0$, and the third has zero mean because
$\int_{\T^2}\operatorname{div}_\xi v=0$.  On the plane $z=1/2$,
the velocity is exactly $(v,0)$, so the plane is invariant and the
particle from $a_*$ is $(Y(t),1/2)$ modulo integers.  This proves that
the fluid lift preserves the event already verified.

The lift costs at most one planar derivative of $v$.  Consequently $U$
obeys all the rapid estimates, is zero initially, and has bounded
gradient on each finite time interval.  Put
\begin{equation}\label{eq:lattice-residual}
 g=\partial_tU+(U\cdot\nabla)U-\nu\Delta U.
\end{equation}
This is the residual-realization construction of
Lemma~\ref{lem:realization}: a prescribed smooth divergence-free
velocity from rest solves the forced equation with this force and
pressure zero.  Here $g$ has zero mean, since its time-derivative
and Laplacian terms integrate to zero and
$(U\cdot\nabla)U=\operatorname{div}(U\otimes U)$.
Product differentiation proves every rapid bound for $g$.  More
explicitly, a derivative of spatial order $k$ of the residual uses
planar derivatives through order $k+3$ and one additional time
derivative, all supplied by \eqref{eq:lattice-slot-bound}.

The field $U$ satisfies the periodic classical hypotheses of
Lemma~\ref{lem:realization}. That lemma gives the stated uniqueness
class, the mean-zero pressure normalization, and the unique global
material flow.

\subsection{Effective evaluation and solenoidal forcing}

The formulas specify all instructions without generating the distinguished
execution.  Given a computable time, a rational enclosure selects
finitely many possible pulse slots.  Within a selected slot, enclose
a real lift of the second coordinate in a rational interval of length
at most $\epsilon_n/2$.  Include every lattice center whose bump
support meets that interval.  The number of candidates is uniformly
bounded, and every omitted term vanishes.  Summing these candidates
avoids deciding whether a computable real lies exactly on a cell
boundary.  Each coefficient requires integer arithmetic and one lookup
in the finite transition table.  Only the loading code $c_0$ uses the
input contents; the later fields use the table and the transformed
input length.

The cutoff formulas have effective derivative bounds, and
\eqref{eq:lattice-slot-bound} supplies an effective tail estimate.
For any fixed $J$, an upper bound on $(n+3)^J2^{-n}$ is computable,
for example by checking finitely many terms until its successive ratio
is bounded by a fixed number below one.  Values of the fields, residual,
and every requested derivative can therefore be evaluated to arbitrary
rational error.  If $\nu$ is an arbitrary fixed positive real instead,
the same algebraic formulas hold, with numerical evaluation relative
to that parameter.

Finally apply Lemma~\ref{lem:projection} to $g$. We have proved its
zero spatial mean and effective estimates at every mixed order and
every temporal weight. Thus the mean-zero solution of
$\Delta\phi=\operatorname{div}g$ gives the effective solenoidal force
\[
 f=g-\nabla\phi,\qquad p=-\phi.
\]
Equation~\eqref{eq:projection} specifies its coefficients, and the
weighted conclusion of the lemma gives \eqref{eq:lattice-rapid}.
The pressure sign preserves $-\nabla p+f=g$, hence precisely the
same velocity and particle event. A competitor for $f$ with pressure
$q$ becomes a competitor for $g$ with pressure $q+\phi$, in the same
normalized classical class. This also proves the projected uniqueness
and completes Theorem~\ref{thm:lattice}.

